\section{Completely monotonic degrees of Stirling remainders}\label{sec:stirling}

\subsection{The question}
A function $f\in C^\infty(0,\infty)$ is completely monotonic (CM) if $(-1)^jf^{(j)}\ge0$ for all $j\ge0$. For $f\to0$ at infinity, the \emph{real completely monotonic degree} is $d_{CM}(f)=\sup\{\rho\ge0: x^\rho f(x)\ \text{is CM}\}$. For $n\ge0$ let
\[
R_n(x)=(-1)^n\Big[\ln\Gamma(x)-\big(x-\tfrac12\big)\ln x+x-\tfrac12\ln2\pi-\sum_{k=1}^n\frac{B_{2k}}{2k(2k-1)x^{2k-1}}\Big],\qquad x>0,
\]
the signed Stirling remainders, with Bernoulli numbers $B_j$. Krasniqi~\cite{Kr} proves
\[
\tfrac83<d_{CM}(R_2)<3,\qquad \tfrac{32}9<d_{CM}(-R_2')<4
\]
(his Theorem 1.2), and, for $\delta_{n}:=d_{CM}(R_n)-(2n-2)$, that $\limsup_n(2n-1)\sup_m\delta_{n,m}\le 2\pi^2/J$ with $J=\int_0^{2\pi}\frac{1-\cos u}u\,\dd u$ (Theorem 1.3; here $\delta_{n,m}$ is the analogous excess for the $m$-th derivative, and $\delta_{n,0}=\delta_n$). He adds that ``the exact real degrees and the optimal asymptotic constants remain open.'' We computed the degrees numerically. The results below are numerical, not proved.

\subsection{Method}
For $n\ge2$, $R_n(x)=\int_0^\infty e^{-xt}h_n(t)\,\dd t$ with $h_n(t)=(-1)^nt^{-1}\big[\frac1{e^t-1}-\frac1t+\frac12-\sum_{k\le n}\frac{B_{2k}t^{2k-1}}{(2k)!}\big]\ge0$, and $-R_n'$ has density $t\,h_n(t)$. By Bernstein's theorem and uniqueness of the Laplace transform, $x^\rho f$ is CM if and only if the Riemann--Liouville derivative $D^\rho h$ of the density is nonnegative on $(0,\infty)$ (for $\rho$ below the order of vanishing of $h$ at $0$ there are no boundary terms). We computed $\rho_c=\sup\{\rho: D^\rho h\ge0\}$; at $\rho_c$ the function $D^\rho h$ first touches zero at a double root $t_*\approx11.47$ (for $R_2$) and $\approx8.04$ (for $-R_2'$). Two independent implementations were used (one with Chebyshev interpolation and QUADPACK algebraic-weight quadrature, one with Gauss--Legendre quadrature after a change of variables), and the second implementation additionally checked the result \emph{without} the Laplace representation by computing exact derivatives of $x^\rho f$ (Leibniz rule and polygamma functions at 300 digits) up to order $200$: all derivatives of the right sign for $\rho=\rho_c-0.02$ and $\rho_c-0.003$, and violations for $\rho=\rho_c+0.02$.

\subsection{Results}
\begin{proposition}[numerical]\label{prop:stirling}
To the digits shown (agreement of the two independent implementations to at least 12 digits),
\[
d_{CM}(R_2)=2.8667042552\ldots,\qquad d_{CM}(-R_2')=3.7984354776\ldots,
\]
both inside the intervals proved in~\cite{Kr}.
\end{proposition}
Table~\ref{tab:stirling} gives $d_{CM}(R_n)$ for further $n$ and the normalised excess $(2n-1)\delta_n$. These values come from our first implementation; the entries marked $^\dagger$ were also recomputed by the independent implementation.

\begin{table}[h]
\centering
\caption{$d_{CM}(R_n)$ and $(2n-1)\delta_n$, $\delta_n=d_{CM}(R_n)-(2n-2)$; and $d_{CM}(-R_n')$ for small $n$. ($^\dagger$: recomputed independently.)}\label{tab:stirling}
\begin{tabular}{rllr}
\toprule
$n$ & $d_{CM}(R_n)$ & $(2n-1)\delta_n$ & $d_{CM}(-R_n')$\\
\midrule
2 & 2.8667042552$^\dagger$ & 2.6001 & 3.7984354776$^\dagger$\\
3 & 4.7074638766$^\dagger$ & 3.5373 & 5.6534285885$^\dagger$\\
4 & 6.5993536696 & 4.1955 & 7.5587377578\\
5 & 8.5213225063$^\dagger$ & 4.6919 & 9.4900039902$^\dagger$\\
6 & 10.4617284606 & 5.0790 & 11.4369053276\\
8 & 14.3760290864 & 5.6404 & 15.3593656569\\
10 & 18.3171745147$^\dagger$ & 6.0263 & --\\
14 & 26.2415343007 & 6.5214 & --\\
18 & 34.1950137509 & 6.8255 & --\\
22 & 42.1635160400$^\dagger$ & 7.0312 & --\\
27 & 52.1360471128 & 7.2105 & --\\
\bottomrule
\end{tabular}
\end{table}

\paragraph{Numerical observation on the sharpness of Krasniqi's constant.} The numbers $(2n-1)\delta_n$ increase monotonically with $n$ in the table. Extrapolating the sequence for $n=2,\dots,27$ by least-squares polynomials of degree $4$--$6$ in $1/n$ or $1/(n+\frac12)$ (using the data for $n\ge8,10,12$) gives limits between $8.0967$ and $8.0977$, median $8.0975$, to be compared with Krasniqi's bound $2\pi^2/J=2\pi^2/\mathrm{Cin}(2\pi)=8.097627\ldots$. The data are therefore compatible with $\lim_n(2n-1)\delta_n=2\pi^2/J$, i.e.\ with his upper bound being attained (at $m=0$); the extrapolation is not a proof, and the spread above is a rough indication of its uncertainty (no error bound).

\paragraph{Caveats.} The computations are quadrature computations with the stability checks described, not interval arithmetic; non-negativity of $D^\rho h$ below $\rho_c$ was checked on grids, and the root $\rho_c$ was located as a tangency (double root) of $D^\rho h$. The entries marked $^\dagger$ agree with the independent implementation to at least 10 digits (for $n=22$ the independent value needed many quadrature nodes to converge; its quoted uncertainty is $3\cdot10^{-11}$); the other entries come from one implementation only.
